% Chapter 1 — Introduction (condensed)

\chapter{Introduction}
\label{Chapter1}

\section{Background and Motivation}

Missing data is a pervasive challenge in real-world machine learning. Medical records omit test results for patients who did not present with relevant symptoms; financial surveys suffer from non-response bias; sensor networks produce gaps due to transmission failures. The proportion of missing values commonly ranges from 5\% to 40\% \cite{vanBuuren2018}, and naive strategies for handling missingness can severely compromise downstream analyses.

Three broad strategies are typically employed. \textbf{Complete case analysis} discards any observation with at least one missing value, sacrificing statistical power and introducing selection bias. \textbf{Single imputation} fills gaps with a single estimate (e.g., column means), distorting the variance structure. \textbf{Multiple imputation} generates several plausible complete datasets and combines inferences across them, appropriately accounting for missingness uncertainty.

Among multiple imputation methods, MICE (Multiple Imputation by Chained Equations; van Buuren \& Groothuis-Oudshoorn, 2011) has become the standard \cite{vanBuuren2011}. MICE imputes conditional means, which are optimal for minimising squared error. However, this optimality comes at a structural cost: imputing $\mathbb{E}[X_j \mid X_{-j}]$ suppresses marginal variance. Van Buuren (2018, \S2.6) explicitly notes that regression imputation ``suppresses variance and produces biased statistical inference'' \cite{vanBuuren2018}. For any analysis requiring the correct marginal distribution --- distributional testing, bootstrap inference, synthetic data generation --- MICE produces systematically distorted results.

This motivates a distributional approach to imputation: rather than minimising pointwise squared error, find imputed values whose \emph{distribution} matches the observed distribution of complete cases. This is precisely the objective of MIGA (Missing data Imputation via Genetic Algorithm).

\section{The MIGA Algorithm}

Figueroa-Garc\'ia, Neruda, and Hernandez-P\'erez (2023) proposed MIGA, published in \textit{Information Sciences} \cite{FigueroaGarcia2023}. MIGA frames missing data imputation as an optimisation problem: find the set of imputed values minimising the statistical distance between the complete rows ($X_A$) and the imputed rows ($X_C$). The fitness function is:

\begin{equation}
    F_r = D_r(\tilde{x}_A,\, \tilde{x}_C) + D_r(\tilde{S},\, I) + D_r(b_A,\, b_C)
    \label{eq:fr}
\end{equation}

where $\tilde{S} = S_A^{-1/2} S_C S_A^{-1/2}$ is the relative covariance, $b$ denotes skewness vectors, and $D_r$ is the Minkowski distance of order $r$ (here $r = \infty$, the Chebyshev norm). A genetic algorithm minimises $F_r$ over $l = 200$ individuals, $G$ generations, and $Q$ independent runs.

The original paper provides no public implementation, no comparison against modern baselines (KNN, MICE), and no MNAR evaluation. This thesis addresses all three gaps.

\section{Contributions}

\textbf{C1 --- Open-source reimplementation.} The first public Python implementation of MIGA, reproducing paper-level results on five UCI benchmark datasets (Iris, Wine, Glass, Haberman, Wholesale), with all underdetermined implementation decisions documented.

\textbf{C2 --- Formal $F_r$--RMSE orthogonality theorem.} We prove that the RMSE-minimising imputation must have $F_r > 0$ whenever conditional variance is non-zero (Proposition~\ref{prop:rmse_implies_fr}), and that any $F_r = 0$ imputation must have $\text{RMSE} > 0$ unless the true missing values are drawn from the same distribution as $X_A$ (Proposition~\ref{prop:fr_implies_rmse}). No single imputation can jointly minimise both objectives.

\textbf{C3 --- First MNAR evaluation.} MIGA is evaluated under three MNAR mechanisms (top-quantile, bottom-quantile, tails). Under top-MNAR on Haberman, MIGA achieves its global minimum $F_r$ (0.810) while simultaneously achieving global maximum RMSE (0.384) --- the sharpest empirical confirmation of C2.

\textbf{C4 --- Systematic baseline comparison.} MIGA, MICE, KNN, and Mean imputation are compared on $F_r$ and RMSE across five datasets with Wilcoxon significance tests (10 seeds). MIGA achieves 1.48--1.97$\times$ lower $F_r$ than MICE on Iris and Haberman ($p < 0.001$); MICE achieves 1.4--2.1$\times$ lower RMSE on every dataset.

\textbf{C5 --- Synthetic $p \times \rho$ characterisation.} A controlled grid experiment on synthetic Gaussian data isolates the effect of dimensionality ($p$) and correlation ($\rho$) on MIGA's advantage. The key finding: failure depends jointly on $(p, \rho)$, not $p$ alone. At $\rho = 0.9$, MIGA loses even at $p = 8$.

\textbf{C6 --- IPW-$F_r$ for MNAR bias correction.} Inverse Probability Weighted $F_r$ re-weights $X_A$ via logistic regression propensity scores. When $n/p \geq 20$, IPW reduces $F_r$ by 1--38\% under directional MNAR. When $n/p < 15$ (Wine), the propensity model overfits and IPW increases $F_r$.

\textbf{C7 --- Adaptive Diversity Injection (AD-MIGA).} The fixed 90\% diversity injection rate in MIGA wastes budget late in runs when the elite pool has stabilised. AD-MIGA replaces it with an exponential decay schedule $\alpha(g) = \exp(-\lambda g/G)$, $\lambda=2.0$, starting at full injection and shifting toward exploitation as the population matures. AD-MIGA-E achieves lower $F_r$ than standard MIGA on all 3 tested datasets (Iris, Haberman, Wholesale) and converges approximately 33 generations earlier on average (105 generations on Iris), while the $F_r$--RMSE orthogonality remains intact.

\textbf{C8 --- Universal $F_r$--RMSE orthogonality via GAIN comparison.} The formal theorem (C2) is extended beyond MICE to neural methods. GAIN \cite{Yoon2018} mixes reconstruction loss (pushing toward conditional means) with an adversarial objective (pushing toward distributional fidelity). We prove (Corollary~\ref{cor:gain_orthogonal}) that no finite $\alpha > 0$ allows GAIN to jointly achieve $F_r = 0$ and $\mathrm{RMSE} = 0$. An empirical $\alpha$-sweep ($\alpha \in \{0,1,10,100,1000,10000\}$, all 5 datasets) reveals three failure regimes: monotone-but-off-frontier (Iris, Wine), instability at high $\alpha$ (Haberman), and MICE-domination on both metrics (Glass). The impossibility holds universally, elevating C2 to a \emph{universal impossibility theorem} for imputation.

\textbf{C9 --- Parallel execution of independent GA runs.} The $Q$ runs share no state and are parallelised via \texttt{n\_jobs} using \texttt{ProcessPoolExecutor} with seed-spawned independent RNGs. Benchmarks across all 5 datasets confirm near-linear scaling (1.94--1.98$\times$ at Q=2). At Q=6, runtime falls from $\sim$200s to $\sim$35s per seed.

\section{Research Objectives}

This thesis pursues four primary objectives:

\begin{enumerate}
    \item \textbf{Reproduce MIGA} from scratch, providing the first open-source Python implementation and verifying replication of the paper's results on five benchmark datasets.
    \item \textbf{Formally prove the $F_r$--RMSE orthogonality}, establishing that no single imputation can jointly minimise both distributional distance and pointwise error.
    \item \textbf{Evaluate MIGA under MNAR}, providing the first systematic assessment of MIGA's behaviour when the missingness mechanism is not at random, including an IPW-$F_r$ correction for selection bias.
    \item \textbf{Characterise when MIGA's $F_r$ advantage holds}, via systematic baseline comparison across five datasets and a controlled $p \times \rho$ synthetic experiment isolating the effects of dimensionality and inter-feature correlation.
\end{enumerate}

\section{Thesis Organisation}

Chapter~\ref{Chapter2} reviews missing data mechanisms, classical and modern imputation methods, genetic algorithms, the original MIGA paper, and formally proves the $F_r$--RMSE orthogonality (Section~\ref{sec:orthogonality}). Chapter~\ref{Chapter3} describes the full MIGA reimplementation, documents every ambiguous design decision, and presents the IPW-$F_r$ extension (Section~\ref{sec:ipw_implementation}) and the synthetic dimension methodology (Section~\ref{sec:synthetic_dim_methods}). Chapter~\ref{Chapter4} reports all experimental results: replication, MNAR evaluation, IPW-$F_r$ comparison, baseline comparison with significance tests, variance preservation, downstream evaluation, and the synthetic $p \times \rho$ experiment. Chapter~\ref{Chapter5} discusses the orthogonality findings, scope conditions, limitations, relationship to prior work, and future directions.
